Теперь давайте поговорим про задачу классификации.
Для начала будем говорить про бинарную классификацию на два класса.
Обобщить эту задачу до задачи классификации на $K$ классов не составит большого труда.

\subsection{Постановка задачи классификации}

Мы хотим обучить линейную модель так, чтобы плоскость, которую она задаёт, как можно лучше отделяла объекты
одного класса от другого.

В идеальной ситуации найдётся плоскость, которая разделит классы: положительный окажется с одной стороны от неё, а
отрицательный с другой.
Выборка, для которой это возможно, называется линейно разделимой.
Увы, в реальной жизни такое встречается крайне редко.

Как обучить линейную модель классификации, нам ещё предстоит понять, но уже ясно, что итоговое предсказание можно
будет вычислить по формуле $y = sgn (w,x_i)$

Сконструируем теперь функционал ошибки так, чтобы он вышеперечисленными проблемами (регрессии) не обладал.
Мы хотим минимизировать число ошибок классификатора, то есть
\[ \mathbb{I} (y_i \neq sgn (w,x_i) ) \rightarrow \min_w \]

Домножим обе части на $y_i$ и немного упростим
\[ \mathbb{I} (y_i < sgn (w,x_i) ) \rightarrow \min_w \]

Величина $M = y_i (w, x_i)$ называется отступом (margin) классификатора.
Такая функция потерь называется misclassification loss.
Свойство.

\begin{itemize}
    \item Отступ положителен, когда $sgn y_i = sgn (w,x_i)$, то есть класс угадан верно; при этом чем больше
    отступ, тем больше расстояние от $x_i$ до разделяющей гиперплоскости, то есть «уверенность классификатора»;
    \item отступ отрицателен, когда $sgn y_i \neq sgn (w,x_i)$, то есть класс угадан неверно; при этом чем больше
    по модулю отступ, тем более сокрушительно ошибается классификатор.
\end{itemize}

От каждого из отступов мы вычисляем функцию $F(M) = \mathbb{I} (M < 0)$.
Она кусочно-постоянная, и из-за этого всю сумму невозможно оптимизировать градиентными методами: ведь её производная
равна нулю во всех точках, где она существует.
Но мы можем мажорировать её какой-нибудь более гладкой функцией, и тогда задачу можно будет решить.

Мы остаёмся в домене supervised learning, то есть уже заданные ответы, на заданных объектах.
Множество ответов имеет меньше чем континуальную мощность.
То есть множество классов, которое мы предсказываем $C$ -- конечное множество.
\[ X \in \mathbb{R}^{n \cross p} , Y \in C^n, \abs{C} < \inf \]
Например, бинарное множество классов, которое мы будем рассматривать большую часть лекции: $C =~{-1, 1}$.
То есть будем заниматься бинарной классификацией, а потом уже перейдем к мультиклассовой.
Соответственно, наша цель -- построить такой классификатор, зависящий от данных, который бы приближал вектор ответов
достаточно хорошо: $c(X) = \tilde{Y} \approx Y $.
Что значит достаточно хорошо -- это мы и обсудим

\subsection{Варианты типов меток}

Пусть теперь наши таргеты $y$ кодируют принадлежность к положительному или отрицательному классу, то есть
принадлежность множеству ${-1,1}$ (в этом параграфе договоримся именно так обозначать классы, хотя в жизни вам будут
нередко встречаться и метки ${0,1}$), а $x$ -- по-прежнему векторы из $\mathbb{R}^D$.

${0,1}$ используются даже чаще.
Разнознаковые обычно в контексте SVM (Support Vector Machines) по историческим и математическим причинам

\subsection{Определение линейного классификатора и геометрическая интерпретация}

Допустим у нас есть матрица признаков и вектор весов, как из этого всего, что уже было в линейной регрессии,
составить линейную классификацию?
Надо делить на классы предсказанный ответ, обрубая по какой-то границе, если у нас два класса.
$c(x) = \mathbb{I} (x \geq 0)$
Здесь граница поставлена ровно 0, почему это не сужает множество классификаторов, которые мы получим и
не нарушает общности?
Ещё на прошлой лекции мы ввели свободный член (bias term) и не пишем его в формулах, чтобы их сократить,
но он тренируется при обучении модели.
Запишем чуть более математично, выражение эквивалентно следующему: $c(x) = sgn (f(x)) = sgn(x^T w)$

Существует геометрическая интерпретация данного выражения: вектор $w$ является нормалью к гиперплоскости в $p$-мерном
пространстве фичей (при равенстве скалярного произведения нулю $x^T w = 0$), с одной стороны гиперплоскости будут одного
класса объекты, с другой -- другого класса объекта.

\subsection{Маржа, определение и ключевое свойство}

Введём понятие margin (отступ): $M_i = y_i f(x_i) = y_i x_i^T w$
Это некоторый аналог ориентированного расстояния, причём если $M_i \leq 0$, значит объект классифицирован неверно.
Отступ обычно используют в теоретических доказательствах (величина проекции на верный класс).

Тогда логично предложить следующую функцию потерь: $Loss = \sum_i M_i \leq 0$
Но у такой loss function будут проблемы, так как это не гладкая функция.
Это главная причина почему мы не будем использовать такой loss, по крайней мере в наших текущих задачах.
Есть ещё один менее очевидный минус: такая метрика ничего нам не говорит об уверенности классификатора.
Представьте объекты как точки в пространстве признаков, есть гиперплоскость классификатора, и было бы разумно
полагать, что объекты на границе, где перемешиваются два класса, будут с меньшей уверенностью разделены
классификатором, чем лежащие далеко \("\)в толще\("\)класса, которые с большей точностью принадлежат своему классу.

\subsection{Выбор весов}

\underline{Общая задача оптимизации}

Веса линейного классификатора выбираются путём минимизации эмпирического риска:

\[ \hat{\bm{w}}, \hat{b} = \arg\min_{\bm{w}, b} \left[ \frac{1}{n} \sum_{i=1}^n \mathcal{L}(y_i, \bm{w}^\top
\bm{x}_i + b) + \lambda R(\bm{w}) \right] \]

где:
\begin{itemize}
    \item $\mathcal{L}(y, z)$ --- функция потерь (loss function)
    \item $R(\bm{w})$ --- регуляризатор
    \item $\lambda \geq 0$ --- параметр регуляризации
\end{itemize}

\underline{Основные подходы к выбору весов}

\begin{enumerate}
    \item \textbf{Перцептрон Розенблатта:} Итеративный алгоритм, сходится для линейно разделимых данных
    \[ \bm{w}^{(t+1)} = \bm{w}^{(t)} + \eta \cdot y_i \bm{x}_i \cdot \mathbb{I}[y_i(\bm{w}^{(t)\top} \bm{x}_i)
    \leq 0] \]
    \item \textbf{Логистическая регрессия:} максимизация правдоподобия с логистической функцией потерь
    \item \textbf{Метод опорных векторов (SVM):} максимизация запаса с hinge-функцией потерь
    \item \textbf{Метод наименьших квадратов:} минимизация MSE (редко используется для классификации)
\end{enumerate}

\subsection{Почему для классификации много разных функций потерь}

Вариантов приблизить ступеньку много и в зависимости от задачи и свойств берут разное приближение